\scanpage{35}
\deleted{On suppose la trace continue.}

\uline{Astuce} : On plonge \unsure{un|une} \unsure{anneau|algèbre} avec une trace non continue, dans un \unsure{anneau} avec une trace continue.

Si \deleted{\unsure{$E_{\varphi}$} \unsure{$\int_{T(A)}$} \unsure{$(t)=$}}

\uline{Prop} 1
\[
\varphi_A(t+0)=\mathop{\rm Inf}\limits_{\substack{E\in P\\\mathrm{Tr}\,E<t}}
\mathop{\rm Sup}\limits_{x\in(1\unsure{+|-}E)H}(Ax,x\unsure{]|)}
\qquad\text{(trace continue)}
\]

% margin note written sideways
\margin{\unsure{:} \unsure{preuve|prouver} \deleted{$\unsure{\Delta_t}(1+A)\unsure{=|\leq}\Delta_t(1+\lvert A\rvert)$} \inserted{$\varphi_{\lvert1+A\rvert}\leq \varphi_{1+\lvert A\rvert}$}} On a en effet, comme on voit aussitôt, (décomposition spectrale) $\varphi_A(t)\geq$ second mb \deleted{\unsure{en}}, Il suffit donc de prouver $\mathop{\rm Sup}\limits_{x\in(1-E)H}(Ax,x)\geq\varphi^-_A(t)$ \inserted{si $\mathrm{Tr}\,E<t$} On est ramené au cas où $A$ est égal à $\lambda F$, $F$ projecteur de trace $\unsure{\geq|>}t$, donc \unsure{à|si} $t$ \unsure{près|fini} \unsure{$=$}. \deleted{\unsure{car ceci}} \inserted{(car $A\geq\varphi_A^{\unsure{-}}(t)F$ avec $\mathrm{Tr}\,F\geq t$)} Il suffit de prouver : Si \uline{$\mathrm{Tr}\,E<\mathrm{Tr}\,F$, alors $(1-E)H\cap FH\neq(0)$}.

En effet, on a \deleted{\unsure{$E$}}, pour
\[
\begin{aligned}
\mathcal X&=E\mathcal H,\ \mathcal Y=F\mathcal H,\\
\mathcal X'&=\mathcal H\ominus\mathcal X,\ \mathcal Y'=\unsure{\mathcal H|\mathcal X}\ominus\mathcal Y :
\end{aligned}
\]
$\mathcal X\ominus(\mathcal X\cap\mathcal Y')\sim\mathcal Y\ominus(\mathcal Y\cap\mathcal X')$ donc $\mathcal Y\cap\mathcal X'=0$ impliquerait \deleted{$\mathcal X\unsure{\sim|=}\unsure{0|\ominus}$} $\mathcal Y\unsure{\prec|\lesssim}\mathcal X$ d'où $\mathrm{Tr}\,F\leq\mathrm{Tr}\,E$, absurde.

\uline{Corollaire} 1 $A$ quelconque implique
\[
\varphi_{\lvert A\rvert}(t+0)=\mathop{\rm Inf}\limits_{\substack{E\in P\\\mathrm{Tr}\,E<t}}
\mathop{\rm Sup}\limits_{x\in(1-E)H}\lVert Ax\rVert\qquad\text{(trace continue)}
\]

\uline{Corollaire 2}
\[
\begin{aligned}
\varphi_{\lvert PA\rvert}\deleted{\unsure{(t)}}&\leq\lVert P\rVert\,\varphi_{\lvert A\rvert}\deleted{\unsure{(t)}}\qquad\text{(trace quelconque),}\\
\varphi_{\lvert AQ\rvert}&\leq\lVert Q\rVert\,\varphi_{\lvert A\rvert}\qquad\deleted{\unsure{\varphi_{\lvert1+A\rvert}\leq\varphi_{1+\lvert A\rvert}}}
\end{aligned}
\]

\uline{Corollaire 3} $\varphi_A$ croît avec \unsure{l'hermitien|l'élément} positif \deleted{\unsure{$A$ positif}} \unsure{$A$}

% margin note written sideways
\margin{\uline{Corollaire 4} $\varphi_{\lvert1+A\rvert}\leq\varphi_{1+\lvert A\rvert}=1+\varphi_{\lvert A\rvert}$}
% slanted
\margin{$\varphi_{\alpha(A)}=\alpha\varphi_A$} \uline{Corollaire 5} Équicontinuité \inserted{des fonctions $A\to\varphi_{\lvert A\rvert}(t)$} \deleted{\unsure{des} $A$ $\varphi_{\lvert A\rvert}(t\unsure{+0})$ en fonction de $A$.}

\uline{Prop} 2 $\varphi_{\lvert A\rvert}=\varphi_{\lvert A^*\rvert}$ Il suffit de le prouver \deleted{que} $\varphi_{AA^*}=\varphi_{A^*A}$ et de le prouver pour $A\in\unsure{a|\mathfrak a}$. \inserted{Posons} \deleted{$A\in\underline a$, \unsure{et} \illegible{1} le prouver} $\unsure{K|X}_{\lvert\unsure{K|A}\rvert A} = K_{\lvert A\rvert}\unsure{\theta\theta|\theta^*}$ ou \unsure{$\int$} \deleted{\unsure{$x\,d.K_{\lvert A\rvert}(xp)|x\,d.K_{\lvert A\rvert}(x,p)$}}
% archive sheet number: 36

